\chapter{EIRENE Continuous Integration}
\label{sec5}

The EIRENE Continuous Integration (CI) pipeline ensures developments to EIRENE do not adversely impact the existing EIRENE code base and allows developers to verify their code across a wide range of EIRENE use cases. The pipeline runs through a series of tests and requires that all the tests pass for successful completion. These tests run on a Gitlab Runner \url{https://docs.gitlab.com/runner/} instance that is triggered whenever code is pushed to the repository, they can also be run manually at any time. The set of tests is defined and controlled by the {.gitlab-ci.yml} file in the EIRENE repository root directory


\section{Overview}
The CI pipeline is broken into the following stages. Each stage contains a number of jobs, the stages run in order, however the \textbf{needs} keyword is used in many of the jobs to allow jobs in later stages to start before the previous stage has finished as long as the \textbf{needs} are met.

\begin{enumerate}
\item preparation
\item build:executable
\item run
\item run\_mpi
\item test
\item preparation:openmp
\item build\:openmp\_executable
\item run\_openmp
\item test\_openmp
\item coverage
\item tag
\item deploy
\item on-schedule
\end{enumerate}

\subsection{Pipeline control}
It is possible to exclude some stages by setting the {EXCLUDE} variables to {true} in order to save time during development. When pushing to either the {develop} or {master} branches the {EXCLUDE} variables are overidden in the {workflow} section so that all stages will run for these branches. The {rules} keyword is used in the {workflow} and also  in the individual jobs to enforce this behaviour. At present the available variables are {EXCLUDE\_SERIAL}, {EXCLUDE\_MPI}, {EXCLUDE\_EMC3}, {EXCLUDE\_OPENMP}. 

The \textbf{rules} keyword allows conditional control of pipeline jobs, each rule is evaluated in order and stops when the first true condition is reached. Detailed documentation of this and other CI commands can be found at \url{https://docs.gitlab.com/ee/ci/yaml/gitlab_ci_yaml.html}.

\section{Continuous Integration Stages}

\subsection{preparation}
The preparation stage contains 2 jobs, {build\_library} and {prepare\_testcases}. The build library stage creates a {buildRelease} directory then calls cmake and make to build the EIRENE library before uploading the {buildRelease} and {libRelease} directories as artefacts. The serial test cases then link to {libRelease/libEIRENE.a}.

The prepare testcases job retrieves the set of EIRENE test cases from the {EIRENE-sample-cases} repository, also on the jugit server. Rather than using git clone a new temporary git repository is created and the remote origin set to the url of the repository. {git fetch} with a depth of 1 is then called using the commit hash in the variable {SAMPLE\_CASE\_VERSION}. When changing the sample cases used in the CI it is this variable that must be changed. {FETCH\_HEAD} is then checked out, resulting in a shallow checkout of depth one, avoiding the download of the complete repository. and saving significant time in the CI pipeline. The {EIRENE-sample-cases} directory is then uploaded as an artefact wit the {.git} directory excluded.

\subsection{build:executable}
In this stage executables are built for each test case. {build\_executable} is an anchor used by the job specification for each sample case. The majority of cases build a standalone executable that is linked to {libRelease/libEIRENE.a}. However for the ITER cases the jobs use the {build\_iter\_executable} anchor and the executable is built directly from the EIRENE source. The build is started from within the makefiles of the ITER cases, which call cmake with the following switches  {-DEIRENE\_INTERFACE=SOLPS-ITER}, {-DEIRENE\_USER-ROUTINES=iter}, and {-DEIRENE\_POSTFIX=\_iter}. An EIRENE executable is then created in binRelease and called from the ITER sample case directories.

\subsection{run}
The run stage executes all of the sample cases with the command {make output}, which ensures the case has been compiled and runs the executable to produce output. Output is generated in both fixed format and json format and stored in the directories {output\_fixed} and {output\_json}. The {run} anchor is used for each of the sample cases.

\subsection{run\_mpi}
The MPI enabled jobs run in a separate stage that runs each tests sequentially due to local resource limits although this will likely be reviewed. The MPI tests are all based on EMC3 coupled cases that are stored in the serial sample repository. This testing is very limited at present and will likely be extended in the future.

\subsection{test}
The test stage uses system {diff} to compare the output of all of the sample cases with reference output stored in each case with the command {make test}. This means that the output must match exactly, so that running the tests on differenct machines or operating systems will almost certainly cause the tests to fail. It is therefore necessary to update the sample cases \cref{loc:sampleCases} when a new Docker image \cref{loc:DockerImage} is used.


\subsection{OpenMP stages}
The OpenMP stages duplicate the serial stages with {\_openmp} appended to the name.  They run after the completion of the serial and MPI tests and compile the tests with OpenMP support using {-DOPENMP=ON} in the cmake command and use 4 OpenMP threads at execution time. Only a subset of testcases are duplicated. A separate branch of the {EIRENE-sample-cases} repository is used for the tests, at the time of writing this is {eirene\_unified\_openmp\_bullseye}. This branch was created from the initial commit, so they can be merged into the original branch containing the serial tests at any time.

The input files for OpenMP are largely the same as for the serial cases apart from the use of the variables {NLIDENT} and {NPTSDEL}. {NLIDENT} is a logical in the first operating mode block of the input file described in \cref{sec2.1} and sets the random seed used by each process or thread to the same value, so that the output can be compared. {NPTSDEL} is not currently described in the input section of the EIRENE manual. It is the sixth entry in the second line of each stratum distribution definition, \cref{sec2.7} in the EIRENE manual, section 7 of the input file. This integer parameter sets the number of particles after which the random seed is reset to the value given by {NINITL}. This means that as long as $NPTS/NPTSDEL$ is an integer product of the number of threads used the serial and threaded cases will use the same initial seed for each block of particles, allowing comparison


\subsubsection{OpenMP comparison}
When comparing output using OpenMP the non-associative nature of multi threaded floating point addition means that small differences are likely to occur in output when compared to serial runs. Because of this it is necessary to compare the relative difference of output and set thresholds on these differences

The {compare\_eirene.py} script compares EIRENE output in 2 directories. \cref{tab:compare_eirene} shows the available command line options that allow control over the output data to be read as well as levels of tolerance and reporting. When called with no arguments the script compares file in the the directories output and output\_ref giving a report for each file where numerical differences exist. The source can be found in any of the openmp specific test case directories, see \cref{loc:sampleCases}.

\begin{table}
  \begin{center}
    \begin{tabular}{ | p{4.2cm} | p{9cm} | }
      \hline
      Option  &  Description \\
      \hline
     -h, --help                  & Show the help message and exit                                \\	
      -d DIR1, --dir1 DIR1        & First directory containing output files (default=output)      \\
      -r DIR2, --dir2 DIR2        & Second directory containing output files (default=output\_ref) \\
      -f FILE, --file FILE        & If specified only this file will be checked                   \\
      -q, --quiet                 & Only prints a message if differences exist                    \\
      -b, --binary                & Check binary files                                            \\
      -v, --verbose               & Prints more detail about the comparisons being made           \\
      -t, --statwarn              & Prints out values of all mismatches above machine precision   \\ 
      -w, --warnings              & Prints out values of all mismatches above statistical error   \\ 
      -e, --errall                & Prints out all mismatches                                     \\
      -s, --single\_precision     & Sets tolerance to single precision                            \\
      -i, --ignore\_small\_values & Ignore values that are much smaller than the mean             \\
      -x, --write\_texfile        & Write differences to a file in tex tabular format             \\
      --ignore\_factor \linebreak IGNORE\_FACTOR  & Set the ratio to the mean at which small values are ignored (default=1000) \\
      \hline
    \end{tabular}
    \caption{Command line options for the compare\_eirene.py script.\label{tab:compare_eirene}}
  \end{center}
\end{table}

The script compares the relative difference between all numeric values in the output files to a series of thresholds, see \cref{tab:compareThresh}. The number of differences at each of level is then reported. The tolerances cannot be changed on the command line and must be changed by hand in the script. The level chosen for Machine Precision is approximate and choice of levels above machine precision is arbitrary however provides a useful representation of differences between output.

\begin{table}
  \begin{center}
    \begin{tabular}{ | p{4cm} | p{4cm} | p{4cm} |  }
      \hline
      Reporting  &  Double Precision  & Single Precision \\
      Level      &  Threshold (Default) & Threshold \\
      \hline
      Machine Precision       &  1$^{-15}$  &  1$^{-7}$ \\
      Statistical Difference  &  1$^{-14}$  &  1$^{-6}$ \\
      Warning                 &  1$^{-12}$  &  1$^{-5}$ \\
      Error                   &  1$^{-3}$   &  1$^{-3}$ \\
      \hline
    \end{tabular}
    \caption{Thresholds and labels for comparison of the relative difference of EIRENE output in the script compare\_eirene.py \label{tab:compareThresh}}
  \end{center}
\end{table}
      
\subsubsection{Exclusion of selected data points}
Some of the EIRENE tallies result in zero or very close to zero values after the addition and subtraction of values that are expected to be the identical due to particles changing species in the relevant cell. In practice machine precision limits mean that the 2 values can differ by a small amount and when one is subtracted from the other a remainder is left. When running serially this remainder is in general the same for repeated runs on the same system, as in the CI. However when OpenMP threading is enabled these values can vary due to the change in ordering of the operations, causing problems for the CI testing as the values can be different by up to a factor 2, causing the pipeline to fail when testing these values.

There is no simple solution available to this problem although in practice only data in the {fort.10} output file in the {ITER\_1573\_openmp} and {2D-D\_slab\_openmp} cases are affected. In order to address this the data points that are affected are excluded from testing by adding a list of points to exclude to the {compare\_eirene.py} script in these directories. The points are selected using the errors generated by in previous runs. Not all of the excluded tallies have been confirmed to be due to machine precision, the tallies that have been identified are represented by the pointers {PMML, PAML} and {PAAT}.

It is likely that changes to the code will result in different cells being affected by this problem, in which case the cells will need to be identified and added to the lists for exclusions. However care must be taken to ensure that this process is in fact the cause of the new errors.

\subsection{coverage}
The coverage stage only runs when pushing to the develop and master branches and can only be used with GNU compilers. This stage clones EIRENE and the sample cases and compiles with the -ftest-coverage flag to produce a coverage report using the Makefile in the EIRENE-sample-cases root dorectory. The GNU tool gcov is used to asses the level of code coverage provided by the serial sample cases and this is reported through the CI. As of the latest release (MsV-2023Mar-PBoerner-beta) the coverage report fails due to a compiler error, this is being investigated.

\subsection{tag, deploy and on-schedule}
There are 3 stages that use the {.tag\_and\_deploy.yml} and {.update\_fork\_branches.yml} files in the EIRENE root directory. These stages do not have any impact on the normal running of the EIRENE CI so are not discussed here.


\section{Sample cases}
\label{loc:sampleCases}
The EIRENE CI uses sample cases from the {EIRENE\_sample-cases} repository at

\url{https://jugit.fz-juelich.de/eirene/EIRENE-sample-cases}

At present two different branches are used for the tests, one for the serial tests and one for the OpenMP tests. This allows development of either without impacting the other. The branch used for serial tests at the time of writing is {eirene\_unified\_bullseye} which contains 24 tests cases covering different geometries and cases. The {eirene\_unified\_openmp\_bullseye} branch is used for the OpenMP cases and contains a subset of 7 of the serial test cases. If either of the sample cases are changed it is necessary to change the relevant git reference in the {.gitlab-ci.yml} file in the root directory of the EIRENE repository.


\section{Development of the Continuous Integration pipeline}


\subsection{Linting and validation}
Changes to the made to the {gitlab-ci.yml} file can be validated on the EIRENE gitlab repository web pages. In the Continuous Integration section there is a button in the top right corner labelled CI lint, or alternatively the web page can be accessed directly at

\url{https://jugit.fz-juelich.de/eirene/eirene/-/ci/lint}

The contents of the gitlab-ci.yml file can be pasted into this page which will validates the syntax and report any errors. In the case of valid syntax it will produce a table of the commands launched by the CI. For developers with permission to write to the develop or master branch a full CI simulation can also be run by selecting the check box \textbf{Simulate pipeline creation for the default branch}. 

\subsection{Docker Images}
\label{loc:DockerImage}
The Docker images used by the EIRENE CI are stored in the repository

\url{git@jugit.fz-juelich.de:eirene/eirene-ci-image.git}

This image can be built locally and the test cases run in a docker environment for testing purposes.

\subsubsection{Using the Docker image}
The use of the docker image on Juelich systems is described in the EIRENE wiki

\url{https://jugit.fz-juelich.de/eirene/eirene/-/wikis/EIRENE Continuous Integration}

In order to run the tests cases on a local Linux machine follow the steps described here. It is assumed that docker is installed on the machine and that you have read access to the following repositories.

\begin{itemize}
\item Docker image \url{git@jugit.fz-juelich.de:eirene/eirene-ci-image.git}
\item Sample cases \url{git@jugit.fz-juelich.de:eirene/EIRENE-sample-cases.git}
\item EIRENE \url{git@jugit.fz-juelich.de:eirene/eirene.git}
\end{itemize}

Create a directory within which one can run the docker container

First build the docker image locally by first cloning the Docker image repository. This will create a directory eirene-ci-image, go into this directory and build the image. 

\begin{lstlisting}[language=bash]
  $ git clone git@jugit.fz-juelich.de:eirene/eirene-ci-image.git
  $ cd eirene-ci-image
  $ sudo docker image build . -t debianbullseye:test
\end{lstlisting}
  
The {debianbullseye:test} tags can be changed to any tags you prefer. This may take some time. It is also necessary to clone the EIRENE and EIRENE-sample-cases repositories. This can be done from either the docker container or from your default environment, if you have ssh keys enabled then this may be easier,  however there is the potential for conflict between the git version on your system and that used in the docker container. You can run the docker image from a directory of your choice, if cloning the EIRENE repository using your local system it should be cloned into the directory that you will run docker from. The sample cases repository should be cloned into the {eirene} directory containing the EIRENE source. Assuming you have already cloned the EIRENE and sample cases repositories you should run the following command under sudo.

\begin{lstlisting}[language=bash]
  $ sudo docker run --privileged=true -it -v `pwd`:`pwd` \
    -w `pwd` debianbullseye:test
\end{lstlisting}

If you have not yet cloned the EIRENE and the sample cases it will now need to be done in the running container. Once this is done you can change directory to the appropriate case and run the {make} command to build the appropriate executable. The {make output} and {make test} will run and test the code. It is unlikely that the tests will pass due to differences in the architecture of the local machine. In this case it is necessary to copy the output generated by a first run into the {output\_ref} directory and then carry out the tests.

